\section*{Introduction}
\addcontentsline{toc}{section}{Introduction}

% TODO: Rédiger l'introduction du chapitre 3
% Points à couvrir:
% 1. Objectif du chapitre (Conception)
% 2. Approche de conception
% 3. Plan du chapitre

[CONTENU À REMPLIR: Introduction du chapitre 3 - Conception]

\section{Architecture générale du système}

% TODO: Décrire l'architecture globale
% Points à inclure:
% - Vue d'ensemble architecturale
% - Composants principaux
% - Technologies utilisées
% - Patterns de conception

\subsection{Architecture matérielle et logicielle}

[CONTENU À REMPLIR: Description de l'architecture matérielle et logicielle]

% Exemple:
% \begin{figure}[H]
%     \centering
%     \includegraphics[width=0.9\textwidth]{images/architecture.png}
%     \caption{Architecture générale du système}
%     \label{fig:architecture}
% \end{figure}

\subsection{Choix technologiques}

% TODO: Justifier les choix technologiques
% Points à inclure:
% - Langage de programmation
% - Frameworks et bibliothèques
% - Base de données
% - Serveur d'application
% - Infrastructure de déploiement

[CONTENU À REMPLIR: Justification des choix technologiques]

% Exemple:
% \begin{itemize}
%     \item \textbf{Frontend:} [Technologie] - [Justification]
%     \item \textbf{Backend:} [Technologie] - [Justification]
%     \item \textbf{Base de données:} [Technologie] - [Justification]
%     \item \textbf{Infrastructure:} [Technologie] - [Justification]
% \end{itemize}

\section{Conception de l'architecture applicative}

% TODO: Détailler l'architecture applicative
% Points à inclure:
% - Couches de l'application
% - Modules et composants
% - Interaction entre modules
% - Patterns utilisés (MVC, MVVM, etc.)

\subsection{Diagrammes de déploiement}

[CONTENU À REMPLIR: Description du diagramme de déploiement]

% Exemple:
% \begin{figure}[H]
%     \centering
%     \includegraphics[width=0.9\textwidth]{images/deployment_diagram.png}
%     \caption{Diagramme de déploiement}
%     \label{fig:deployment}
% \end{figure}

\subsection{Diagrammes de séquence}

% TODO: Créer des diagrammes de séquence pour les scénarios clés
% Points à inclure:
% - Interactions entre objets
% - Flux de messages
% - Ordre temporel des actions

[CONTENU À REMPLIR: Description des diagrammes de séquence]

% Exemple:
% \begin{figure}[H]
%     \centering
%     \includegraphics[width=0.9\textwidth]{images/sequence_diagram.png}
%     \caption{Diagramme de séquence: [Scenario name]}
%     \label{fig:sequence_diagram}
% \end{figure}

\section{Conception de la base de données}

% TODO: Concevoir la base de données
% Points à inclure:
% - Schéma de la base de données
% - Tables et colonnes
% - Contraintes et indices
% - Normalization

\subsection{Schéma physique de la base de données}

[CONTENU À REMPLIR: Description du schéma physique de la base de données]

% Exemple:
% \begin{figure}[H]
%     \centering
%     \includegraphics[width=0.9\textwidth]{images/database_schema.png}
%     \caption{Schéma physique de la base de données}
%     \label{fig:db_schema}
% \end{figure}

\subsection{Scripts de création de la base de données}

[CONTENU À REMPLIR: Informations sur les scripts de création]

% Exemple:
% \begin{verbatim}
% CREATE TABLE users (
%     user_id INT PRIMARY KEY AUTO_INCREMENT,
%     username VARCHAR(255) NOT NULL UNIQUE,
%     email VARCHAR(255) NOT NULL UNIQUE,
%     password_hash VARCHAR(255) NOT NULL,
%     created_at TIMESTAMP DEFAULT CURRENT_TIMESTAMP
% );
% \end{verbatim}

\section{Design des interfaces utilisateur}

% TODO: Concevoir les interfaces utilisateur
% Points à inclure:
% - Mockups/wireframes
% - Design patterns
% - Principes UX/UI
% - Navigation

\subsection{Maquettes et wireframes}

[CONTENU À REMPLIR: Description des maquettes et wireframes]

% Exemple:
% \begin{figure}[H]
%     \centering
%     \includegraphics[width=0.8\textwidth]{images/wireframe.png}
%     \caption{Wireframe de la page d'accueil}
%     \label{fig:wireframe}
% \end{figure}

\subsection{Prototypes interactifs}

[CONTENU À REMPLIR: Description des prototypes interactifs]

\section{Sécurité et performance}

% TODO: Planifier les aspects sécurité et performance
% Points à inclure:
% - Authentification et autorisation
% - Chiffrement des données
% - Validation des entrées
% - Optimisation des requêtes
% - Caching

\subsection{Stratégie de sécurité}

[CONTENU À REMPLIR: Description de la stratégie de sécurité]

% Exemple:
% \begin{itemize}
%     \item Authentification: [Méthode]
%     \item Autorisation: [Approche]
%     \item Chiffrement: [Technique]
%     \item Validation: [Méthode]
% \end{itemize}

\subsection{Optimisations de performance}

[CONTENU À REMPLIR: Description des optimisations de performance]

\section*{Conclusion}
\addcontentsline{toc}{section}{Conclusion}

% TODO: Conclure le chapitre 3
% Points à inclure:
% - Synthèse du design architectural
% - Avantages de l'architecture proposée
% - Transition vers la réalisation

[CONTENU À REMPLIR: Conclusion du chapitre 3]
